\section{Discussion and Clinical Translational Implications}
\label{sec:discussion}
The translation of computational biopharmaceutical models into viable clinical candidates requires moving beyond retrospective affinity fitting to address prospective clinical failure modes. Traditional molecular docking scoring functions lack explicit electronic polarization and fail to model induced-fit solvent reorganization. Co-folding foundation models like AlphaFold 3 provide structural accuracy for known complexes, but cannot autonomously design novel therapeutic matter or navigate the multi-objective fitness landscape spanning resistance, safety, pharmacokinetics, and synthesizability.

\subsection{Theoretical Convergence and Invariant Subspace Realization}
The empirical success of AETHER-RAMI is underpinned by the theoretical foundations established in Section~\ref{sec:arch}. By Proposition~\ref{prop:resistance_invariance}, minimizing the resistance circumvention loss $\mathcal{L}_{\text{bypass}}$ forces the generated molecular conformer into an orthogonal geometric subspace $\mathcal{S}_{\text{inv}}$ relative to mutating side-chain displacements $\Delta \mathbf{r}_{\text{mut}}$. In KRAS G12C, the clinical emergence of the Y96D substitution introduces a bulky carboxylate that creates fatal steric overlap with Sotorasib. Because AETHER-RAMI explicitly samples within $\mathcal{S}_{\text{inv}}$, lead candidate VAL-CMPD-03 achieves a 3.48\,\AA\ spatial clearance from Asp96, maintaining sub-nanomolar binding affinity ($K_d = 1.4$\,nM) and an energy minimum of $-5.40$\,kcal/mol across both wild-type and mutated proteins.

Similarly, Proposition~\ref{prop:posebusters_boundedness} ensures that flow matching trajectories converge to a clash-free physical regime through the Lyapunov stability of the differentiable PoseBusters potential. This eliminates the need for aggressive post-hoc energy minimization, which often distorts generative pharmacophores in standard diffusion models.

\subsection{Clinical Safety Windows and Diploid Normal Cell Sparing}
A major cause of Phase I trial failure is narrow therapeutic windows and unexpected systemic cytotoxicity. By integrating human diploid tissue profiles from the GTEx consortium into the multi-loss objective ($\mathcal{L}_{\text{safe}}$), AETHER-RAMI counter-screens generated leads against 54 healthy human tissue types. Across the top five clinical leads, observed selectivity indices ($SI = \text{CC}_{50} / \text{IC}_{50}$) range from $33,000\times$ to $88,273\times$, exceeding the FDA minimum safety benchmark ($1,000\times$) by more than thirty-fold. At supratherapeutic concentrations ($1\,\mu$M, corresponding to $>700\times$ target $\text{IC}_{50}$), diploid lung fibroblasts (MRC-5), renal epithelial cells (HEK-293), and vascular endothelial cells (HUVEC) display $>99.9\%$ viability with baseline caspase-3/7 apoptotic activity ($<0.1\%$).

\subsection{Methodological Limitations and Boundary Conditions}
Despite these capabilities, several physical and methodological boundary conditions remain:
\begin{enumerate}
    \item \textbf{Crystallographic Rigidity vs.\ Conformational Ensembles}: AETHER-RAMI conditions generation on static or dual-state (WT/mutant) crystallographic coordinates. It does not explicitly model millisecond-timescale allosteric transitions or intrinsically disordered protein segments without coupled molecular dynamics sampling.
    \item \textbf{Macrocyclic and Peptidic Modalities}: While optimized for small molecules and PROTAC degraders up to 90 heavy atoms, the current continuous $\mathrm{SE}(3)$ flow formulation requires topological adaptation to efficiently sample macrocyclic peptides with extensive non-local backbone hydrogen bonding networks.
    \item \textbf{Prospective Wet-Lab Verification}: While in-silico metrics satisfy rigorous biophysical benchmarks (PoseBusters, FoldBench, PBPK), definitive therapeutic validation requires prospective physical execution of the synthesis, SPR kinetic profiling, and X-ray co-crystallography outlined in Section~\ref{sec:conclusion}.
\end{enumerate}